\documentclass[11pt,letterpaper]{article}
\usepackage[T1]{fontenc}
\usepackage[utf8]{inputenc}
\usepackage[margin=0.88in,top=0.84in,bottom=0.83in,headheight=13pt,headsep=20pt,footskip=28pt]{geometry}
\usepackage{newpxtext}
\usepackage{amsmath,amsthm}
\usepackage{newpxmath}
\usepackage{microtype}
\usepackage{xcolor}
\definecolor{Forest}{HTML}{30392C}
\definecolor{Olive}{HTML}{687144}
\definecolor{Muted}{HTML}{65685F}
\definecolor{Sage}{HTML}{CBD0BC}
\definecolor{Pale}{HTML}{F2F4EB}
\usepackage{mathtools}
\usepackage{booktabs,tabularx,longtable,array}
\usepackage{enumitem}
\usepackage{titlesec}
\usepackage{fancyhdr}
\usepackage[most]{tcolorbox}
\usepackage{needspace}
\PassOptionsToPackage{obeyspaces}{url}
\usepackage{xurl}
\usepackage[colorlinks=true,linkcolor=Olive,citecolor=Olive,urlcolor=Olive,bookmarksnumbered=true,pdfencoding=auto,psdextra]{hyperref}
\hypersetup{pdftitle={Cover exactingness between strongly compact cardinals},pdfauthor={Prepared for Vladimir Reshetnikov},pdfsubject={A conditional inconsistency theorem using omega-club amenability and stationary seeds},pdfkeywords={large cardinals, cover exacting, HCD, strongly compact, omega-club amenability, stationary seeds, conditional inconsistency}}
\urlstyle{same}
\setlength{\parindent}{0pt}
\setlength{\parskip}{5.1pt plus 1pt minus .4pt}
\linespread{1.035}
\setlength{\emergencystretch}{2em}
\setcounter{tocdepth}{1}
\setcounter{secnumdepth}{2}
\titleformat{\section}{\Large\bfseries\color{Forest}}{\thesection}{.6em}{}
\titleformat{\subsection}{\large\bfseries\color{Forest}}{\thesubsection}{.55em}{}
\titleformat{\subsubsection}{\normalsize\bfseries\color{Forest}}{}{0pt}{}
\titlespacing*{\section}{0pt}{18pt plus 3pt minus 2pt}{8pt}
\titlespacing*{\subsection}{0pt}{12pt plus 2pt minus 1pt}{5pt}
\titlespacing*{\subsubsection}{0pt}{9pt}{3pt}
\setlist[itemize]{leftmargin=1.4em,itemsep=2pt,topsep=3pt,parsep=0pt}
\setlist[enumerate]{leftmargin=1.7em,itemsep=3pt,topsep=4pt,parsep=0pt}
\pagestyle{fancy}
\fancyhf{}
\fancyhead[L]{\sffamily\fontsize{9}{11}\selectfont\color{Muted}LARGE CARDINALS AT THE INCONSISTENCY FRONTIER}
\fancyhead[R]{\sffamily\fontsize{9}{11}\selectfont\color{Muted}CONTINUATION IV}
\fancyfoot[L]{\small\color{Muted}Stationary seeds and a conditional inconsistency}
\fancyfoot[R]{\small\color{Muted}\thepage}
\renewcommand{\headrulewidth}{.35pt}
\renewcommand{\footrulewidth}{0pt}
\fancypagestyle{plain}{\fancyhf{}\fancyfoot[L]{\small\color{Muted}Stationary seeds and a conditional inconsistency}\fancyfoot[R]{\small\color{Muted}\thepage}\renewcommand{\headrulewidth}{0pt}}
\tcbset{colback=Pale,colframe=Sage,colbacktitle=Sage,coltitle=Forest,fonttitle=\bfseries,boxrule=.5pt,arc=3pt,left=9pt,right=9pt,top=6pt,bottom=6pt,toptitle=3pt,bottomtitle=3pt,before skip=8pt,after skip=8pt,breakable}
\newtcolorbox{assessment}[1]{title={#1}}
\theoremstyle{definition}
\newtheorem{definition}{Definition}[section]
\newtheorem{theorem}[definition]{Theorem}
\newtheorem{proposition}[definition]{Proposition}
\newtheorem{lemma}[definition]{Lemma}
\newtheorem{corollary}[definition]{Corollary}
\newtheorem{remark}[definition]{Remark}
\newtheorem{question}[definition]{Question}
\newtheorem{inputthm}{Input}
\renewcommand{\theinputthm}{\Alph{inputthm}}
\numberwithin{equation}{section}
\newcommand{\ZFC}{\mathsf{ZFC}}
\newcommand{\ZF}{\mathsf{ZF}}
\newcommand{\AC}{\mathsf{AC}}
\newcommand{\GA}{\mathsf{GA}}
\newcommand{\HOD}{\mathrm{HOD}}
\newcommand{\OD}{\mathrm{OD}}
\newcommand{\CD}{\mathrm{CD}}
\newcommand{\HCD}{\mathrm{HCD}}
\newcommand{\UF}{\mathrm{UF}}
\newcommand{\Ex}{\operatorname{Ex}}
\newcommand{\UEx}{\operatorname{UEx}}
\newcommand{\Ext}{\operatorname{Ext}}
\newcommand{\SC}{\operatorname{SC}}
\newcommand{\CEx}{\operatorname{CEx}}
\newcommand{\REx}{\operatorname{REx}}
\newcommand{\Con}{\operatorname{Con}}
\newcommand{\crit}{\operatorname{crit}}
\newcommand{\cf}{\operatorname{cf}}
\newcommand{\ot}{\operatorname{ot}}
\newcommand{\ran}{\operatorname{ran}}
\newcommand{\rank}{\operatorname{rank}}
\newcommand{\lcm}{\operatorname{lcm}}
\newcommand{\Ord}{\mathrm{Ord}}
\newcommand{\Pow}{\mathcal P}
\newcommand{\id}{\mathrm{id}}
\newcommand{\restr}{\mathbin{\upharpoonright}}
\newcommand{\Izero}{\mathrm{I0}}
\newcommand{\Itwo}{\mathrm{I2}}
\newcommand{\Ithree}{\mathrm{I3}}
\newcommand{\Iwf}{\mathrm{I3}_{\mathrm{wf}}(0)}
\newcommand{\Sel}{\operatorname{Sel}}
\newcommand{\FF}{\mathcal F}
\newcommand{\PP}{\mathbb P}
\newcommand{\Clam}{\mathcal C_\lambda}
\newcommand{\Rig}{\operatorname{Rig}}
\newcommand{\Spec}{\operatorname{Spec}}
\newcommand{\ch}{\operatorname{ind}}
\newcommand{\CF}{\mathsf{CF}}
\newcommand{\src}[1]{\unskip\ {\normalfont\small\sffamily\color{Olive}[#1]}\ }
\newcommand{\R}[1]{\textsf{R#1}}
\newcolumntype{Y}{>{\raggedright\arraybackslash}X}
\newcolumntype{P}[1]{>{\raggedright\arraybackslash}p{#1}}
\newcolumntype{C}{>{\centering\arraybackslash}p{.034\textwidth}}
\newcommand{\yes}{$\bullet$}
\newcommand{\prt}{$\circ$}
\newcommand{\arxiv}[1]{\href{https://arxiv.org/abs/#1}{arXiv:#1}}
\newcommand{\doi}[1]{\href{https://doi.org/#1}{doi:\nolinkurl{#1}}}


\newcommand{\calF}{\mathcal F}
\newcommand{\calI}{\mathcal I}
\newcommand{\calS}{\mathcal S}
\newcommand{\calU}{\mathcal U}
\newcommand{\calD}{\mathcal D}
\newcommand{\calP}{\mathcal P}
\newcommand{\tc}{\operatorname{tc}}
\newcommand{\Ffine}{\mathsf{F}_{\omega_1}}
\newcommand{\pos}{\mathord{+}}
\newcommand{\st}{\mathrm{st}}
\newcommand{\NS}{\mathrm{NS}}
\newcommand{\card}{\operatorname{card}}
\newcommand{\wca}{\text{$\omega$-club amenable}}
\newcommand{\powN}[2]{(2^{#1})^{#2}}
\newcommand{\New}{\textsf{New deduction}}
\newcommand{\Old}{\textsf{Inherited}}
\begin{document}
\thispagestyle{empty}
{\sffamily\bfseries\small\color{Olive}SET THEORY\quad /\quad RESEARCH CONTINUATION IV}

\vspace{13pt}
{\fontsize{28}{31}\selectfont\bfseries\color{Forest}Cover exactingness between\\ strongly compact cardinals\par}
\vspace{10pt}
{\fontsize{14}{18}\selectfont A conditional inconsistency theorem from\\ $\omega$-club amenability and stationary seeds\par}
\vspace{14pt}
{\color{Muted}Prepared for Vladimir Reshetnikov\hfill 18 September 2026\par}
\vspace{3pt}
{\color{Sage}\hrule height .6pt}
\vspace{12pt}

\textbf{Abstract.}
We continue the completeness-barrier analysis in the supplied
\emph{Large Cardinals Synthesis}. The additional observation is that
$\HCD(\eta)$ contains the trace of every ambient $\omega$-club filter.
A stationary-seed argument then shows that an $\omega$-club-amenable
set-forcing ground must have the $\delta$-cover property whenever
$\delta$ is strongly compact in the extension. Such a ground cannot
regard an ordinal above $\delta$ as regular if it has ambient cofinality
$\omega$. Combining this with the inherited completeness barrier and
Goldberg's ground theorem gives the following conditional inconsistency.

\begin{assessment}{Main result: the two strongly compact cardinals cannot coexist with the cover}
There are no cardinals
\[
\delta<\lambda\le\gamma<\kappa
\]
such that $\delta$ and $\kappa$ are strongly compact and $\lambda$ is
$\gamma$-cover exacting. Equivalently,
\[
\ZFC\vdash
\neg\exists\delta,\lambda,\gamma,\kappa\,
\begin{bmatrix}
\delta<\lambda\le\gamma<\kappa\ \wedge\ \SC(\delta)\\[-2pt]
{}\wedge\ \CEx_\gamma(\lambda)\ \wedge\ \SC(\kappa)
\end{bmatrix}.
\]
No Ground Axiom, stabilization of the $\HCD$ hierarchy, or preservation
of a particular successor cardinal is assumed.
\end{assessment}

We also obtain a local strengthening. If $\delta<\lambda$ is strongly
compact, $\CEx_\gamma(\lambda)$ holds, and a high-completeness model
$N=\HCD(\eta)$ satisfies Choice, then every sufficiently high ambient
regular cardinal is measurable in $N$. More precisely, it suffices that
$\theta>(2^\lambda)^N$. The proof bounds the number of stationary cells
that can survive in an $N$-partition and includes a proof of the
classical filter-atomicity lemma used in this last step.

\textbf{Status and attribution.}
These are new deductions \emph{relative to the supplied synthesis}, not
claims of historical priority. The stationary-seed method and the
underlying HOD dichotomy are due to Goldberg. The cover-exacting
characterizations are due to Blue--Goldberg. Detailed English proofs of
the deductions are supplied; major external inputs are explicitly
separated from them. This report is unrefereed and has not been formally
verified in Lean. It does not refute cover exactingness by itself, nor
does it settle the unrestricted problem of a cover-exacting cardinal
above just one strongly compact cardinal.

\newpage
\tableofcontents
\vspace{9pt}
\begin{assessment}{The argument in one line}
\[
\begin{gathered}
\CEx_\gamma(\lambda)
\ \Longrightarrow\ \cf^{\HCD(\kappa)}(\lambda)=\lambda,
\qquad \cf^V(\lambda)=\omega,\\
\SC(\kappa),\ \kappa>\gamma
\ \Longrightarrow\ \HCD(\kappa)\text{ is a ground},\\
\SC(\delta),\ \delta<\lambda,
\ \omega\text{-club amenability of that ground}
\ \Longrightarrow\ \cf^{\HCD(\kappa)}(\lambda)<\delta.
\end{gathered}
\]
The contradiction is between the first and last lines. The additional
covering in the last line comes from stationary seeds; it is not an
illicit increase in the completeness of a given ultrafilter.
\end{assessment}

\newpage
\section{What is added to the previous research}\label{sec:contribution}

The supplied synthesis \cite{archive}, especially its completeness
barrier and its sections on a strongly compact cardinal above the cover
bound, already proves
\begin{equation}\label{eq:inherited-gap}
\CEx_\gamma(\lambda)
\quad\Longrightarrow\quad
\cf^V(\lambda)=\omega
\quad\text{and}\quad
\cf^{\HCD(\eta)}(\lambda)=\lambda
\quad(\eta\ge q(\gamma)),
\end{equation}
where
\[
q(\gamma)=
\begin{cases}
\gamma,&\gamma\text{ is singular},\\
\gamma^+,&\gamma\text{ is regular}.
\end{cases}
\]
It also observes that an upper strongly compact $\kappa>\gamma$ makes
$\HCD(\kappa)$ a proper ground. With a lower strongly compact
$\delta<\lambda$, its ``two strongly compacts'' corollary gives strict
separation between two complete-definability models. Its later
``collapse or destroy'' theorem leaves a disjunction between collapsing
the old successor and destroying a particular stationary set.

The missing additional ingredient is not a stronger completeness
assumption. It is the fact that the high-completeness model can
\emph{recognize the trace of ambient $\omega$-club positivity}. This
permits one to put a small stationary seed \emph{inside the image of
that model}. A lower strongly compact cardinal then supplies small
covers in the high model, contradicting \eqref{eq:inherited-gap}.

\begin{theorem}[No strongly compact sandwich]\label{thm:main}
In $\ZFC$, suppose $\delta<\lambda\le\gamma$ are cardinals,
$\delta$ is strongly compact, and $\CEx_\gamma(\lambda)$ holds.
Then no strongly compact cardinal is greater than $\gamma$.
\end{theorem}

The proof appears in Section~\ref{sec:mainproof}, after a general
no-ground theorem. A second, shorter verification uses Goldberg's
published covering dichotomy in Section~\ref{sec:crosscheck}.

\begin{assessment}{How the earlier alternatives change}
Under \emph{both} strongly compact hypotheses, the forcing extension
postulated in the earlier analysis does not exist. Thus we do not
choose the collapse branch or the stationary-destruction branch of that
analysis: we rule out their common premise. Without the upper strongly
compact, the original coexistence question is not decided here.
\end{assessment}

A separate new consequence, which does not assume an upper strongly
compact, is the quantitative stationary-width theorem of
Section~\ref{sec:width}. Whenever a high $\HCD$ model satisfies $\ZFC$,
its ambient $\omega$-club filters have no internally indexed disjoint
positive family of length $\lambda$. This yields an explicit threshold
for measurability in that inner model. The Axiom of Choice in the high
model is an actual hypothesis in this consequence; it is not silently
inferred from the lower strongly compact.

\section{Definitions and the external inputs}\label{sec:setup}

\subsection{Ambient and internal conventions}

The ambient universe $V$ satisfies $\ZFC$. An inner model is transitive,
contains all ordinals, and is a class of $V$. Unadorned cardinalities,
ultrafilter completeness, cofinalities, and stationarity are ambient.
A superscript $N$ specifies computation in $N$. In particular,
$P_\delta(\mu)=\{a\subseteq\mu:|a|^V<\delta\}$.

We use only set-parameter-definable inner models in arguments involving
an ultrapower. If $N=\{x:\varphi(x,p)\}$ and $j:V\to M$ is an
ultrapower embedding, the notation $N^j$ means
\[
N^j=\{x\in M:M\models\varphi(x,j(p))\}.
\]
It does \emph{not} mean the generally nontransitive pointwise image
$j``N$. For $N=\HCD^V(\eta)$, it is exactly
$N^j=\HCD^M(j(\eta))$. All uses of $j$ below arise from a set
ultrafilter; the relevant restrictions of $j$ and images of sets are
sets definable from that ultrafilter. No additional axiom asserting a
class embedding is being imposed.

A ground is an inner model $W\models\ZFC$ such that
$V=W[G]$ for a set forcing $\PP\in W$ and a $W$-generic filter $G$.
A \emph{$\delta$-cover} of $x$ in $N$ is a set $a\in N$ with
$x\subseteq a$ and $|a|^N<\delta$.

\begin{lemma}[Comparing smallness]\label{lem:smallness}
If $N\models\ZFC$ is an inner model, $\delta$ is a cardinal of $V$,
and $a\in N$, then $|a|^V<\delta$ if and only if $|a|^N<\delta$.
\end{lemma}
\begin{proof}
The ordinal $\delta$ is also a cardinal in $N$, since a bijection in
$N$ witnessing otherwise would remain a bijection in $V$. An
$N$-bijection from an ordinal below $\delta$ onto $a$ proves ambient
smallness. Conversely, if $N$ regarded $a$ as having size at least
$\delta$, Choice in $N$ would give an injection $\delta\to a$ in $N$,
and hence in $V$. That contradicts $|a|^V<\delta$.
\end{proof}

\subsection{Cover exactingness and complete definability}

For $\alpha>\lambda$ and $Y\subseteq V_\alpha$, a relative
$\lambda$-exacting witness is an elementary embedding
\[
j:(X,\in,X\cap Y)\longrightarrow(V_\alpha,\in,Y)
\]
where the source is an elementary substructure of the target,
$V_\lambda\cup\{\lambda\}\subseteq X$, $j(\lambda)=\lambda$, and
$j\restr\lambda\ne\id$. The domain need not be transitive. Write
$\REx(\lambda,Y)$ when such witnesses exist at all sufficiently large
heights containing the predicate $Y$. The equivalent all-heights
version may also be used.

\begin{definition}[Uniform cover bound]
For infinite cardinals $\lambda\le\gamma$, $\CEx_\gamma(\lambda)$ means
that for every $\alpha>\lambda$ and every $y\in V_\alpha$, some
$Y\subseteq V_\alpha$ satisfies $y\in Y$, $|Y|\le\gamma$, and admits a
relative $\lambda$-exacting witness at height $\alpha$.
\end{definition}

Here $\gamma$ is fixed before the height and the parameter are chosen;
$y\in Y$ must not be replaced by $y\subseteq Y$. Ordinary exactingness
is obtained by forgetting the extra predicate and cover requirement.

For an infinite cardinal $\eta$, $\CD(\eta)$ consists of sets definable
in $V$ from finitely many ordinal parameters and finitely many
$\eta$-complete ultrafilters on ordinals. Completeness means closure
under intersections of fewer than $\eta$ members. Put
\[
\HCD(\eta)=\{x:\tc(\{x\})\subseteq\CD(\eta)\}.
\]
These definitions agree with Goldberg's complete-definability classes
\cite[\S4.1]{unique}. We have
\begin{equation}\label{eq:monotonic}
\eta\le\xi\quad\Longrightarrow\quad
\CD(\xi)\subseteq\CD(\eta),\qquad
\HCD(\xi)\subseteq\HCD(\eta).
\end{equation}
Every ordinal-definable set belongs to $\CD(\eta)$. For a set of
ordinals $a$, membership in $\CD(\eta)$ is equivalent to membership in
$\HCD(\eta)$.

\subsection{Imported theorems, with their exact roles}

\begin{inputthm}[Local Kunen inconsistency]\label{in:kunen}
There is no nontrivial elementary embedding
$V_{\rho+2}\to V_{\rho+2}$ in $\ZFC$ \cite{kunen}.
\end{inputthm}

\begin{inputthm}[Blue--Goldberg]\label{in:BG}
The following consequences of cover exactingness are used
\cite[Lemma~3.3 and Proposition~3.4]{cover}.
If $\REx(\lambda,Y)$ and $z\in\OD_{\{Y\}}$, then
$\REx(\lambda,z)$. Also, if $\CEx_\gamma(\lambda)$ holds, then for all
sufficiently large $\alpha$ the set
\[
\{\sigma\in P_{\gamma^+}(V_\alpha):\REx(\lambda,\sigma)\}
\]
is stationary. In particular it meets the club of elementary hulls
containing any prescribed finite set of parameters.
\end{inputthm}

These are explicitly stated lecture-note inputs, not a newly claimed
consistency construction. The supplied synthesis also gives a
persistent-cover proof of the exact-hull form of this input. We use the
stationary form to keep the inherited portion short, and reprove the
trace argument in Section~\ref{sec:barrier}.

\begin{inputthm}[Goldberg's complete-definability ground theorem]\label{in:ground}
For every infinite $\eta$, $\HCD(\eta)\models\ZF$. If $\kappa$ is
strongly compact, then $\HCD(\kappa)\models\ZFC$ and is a set-forcing
ground of $V$ \cite[Proposition~4.2, Theorems~4.3 and~4.9, arXiv v2]{unique}.
\end{inputthm}

The theorem numbers in the journal version are shifted: the ground
statement is Theorem~4.10. The arXiv v2 statement and proof were checked
directly. Strong compactness, not supercompactness, suffices. The
stronger hypothesis in the shorter 2026 lecture-note summary should not
be substituted for the actual hypothesis of this theorem. No bound on
the size of the ground forcing is used in the main result.

We also use the usual ultrafilter formulation of strong compactness:
for every ordinal $\theta\ge\delta$, a strongly compact $\delta$ admits
a $\delta$-complete fine ultrafilter on $P_\delta(\theta)$. This is the
form needed by the stationary-seed method in
\cite[Proposition~2.3]{dichotomy}. Fineness means that
$\{a:\xi\in a\}$ is in the ultrafilter for every $\xi<\theta$.
Countable completeness makes its ultrapower well founded. These
standard large-cardinal and ultrapower facts are external background;
we do not present them as additional new results.

\section{Reconstructing the inherited completeness barrier}\label{sec:barrier}

This section identifies exactly what is carried over from the supplied
reports. The rest of the paper needs only its final conclusion.

\begin{lemma}[Fixed ordinals of a rank self-embedding]\label{lem:critical}
Let $\lambda$ be an uncountable cardinal and let
$e:V_\lambda\to V_\lambda$ be nontrivial. If
$\kappa_0=\crit(e)$ and $\kappa_{n+1}=e(\kappa_n)$, then
\[
\sup_{n<\omega}\kappa_n=\lambda,
\qquad e(\xi)>\xi\quad(\kappa_0\le\xi<\lambda).
\]
Consequently $\cf^V(\lambda)=\omega$, and every fixed ordinal below
$\lambda$ is below $\crit(e)$.
\end{lemma}
\begin{proof}
Suppose $\rho=\sup_n\kappa_n<\lambda$. The critical sequence and the
sets needed to define its supremum belong to $V_\lambda$: since
$\lambda$ is an uncountable cardinal, adding finitely many ranks, or
$\omega$ ranks, to an ordinal below $\lambda$ still stays below
$\lambda$. Applying $e$ to the sequence shifts it by one place, so
$e(\rho)=\rho$. The restriction of $e$ to $V_{\rho+2}$ is elementary
and nontrivial, contradicting Input~\ref{in:kunen}. Thus the critical
sequence is cofinal in $\lambda$. If
$\kappa_n\le\xi<\kappa_{n+1}$, monotonicity on ordinals gives
$e(\xi)\ge e(\kappa_n)=\kappa_{n+1}>\xi$. This proves both assertions.
\end{proof}

An exacting witness restricts to such an $e$. Indeed $V_\lambda$ is
definable from the fixed parameter $\lambda$, is contained in its
domain, and satisfaction for that rank segment is preserved. This
argument does not assume that the full elementary submodel is
transitive.

\begin{lemma}[A short cofinal set cannot be a preserved predicate]\label{lem:cofinal}
If $a\subseteq\lambda$ is cofinal and $\ot(a)<\lambda$, then
$\REx(\lambda,a)$ fails.
\end{lemma}
\begin{proof}
At a sufficiently large witnessing height, the set $a$ is the unique
set represented by its unary predicate. Thus $a\in X$ and $j(a)=a$.
Its order type $\tau$ and its increasing enumeration $h:\tau\to a$
are definable from $a$, belong to $X$, and are fixed by $j$.
Since $\tau<\lambda$ is fixed, Lemma~\ref{lem:critical} implies
$\tau<\crit(j)$. Every $\xi<\tau$ is therefore fixed, and
\[
j(h(\xi))=j(h)(j(\xi))=h(\xi).
\]
The cofinal set $a$ would consist of fixed ordinals below $\lambda$,
although all those fixed ordinals lie below $\crit(j)$. This is
impossible.
\end{proof}

\begin{lemma}[Reconstruction from a small trace]\label{lem:trace}
Suppose $U$ is an $\eta$-complete ultrafilter on an ordinal $\tau$,
$U,\tau\in\sigma\prec V_\alpha$, the height is sufficiently large,
and $|\sigma|<\eta$. Then $U\in\OD_{\{\sigma\}}$.
\end{lemma}
\begin{proof}
The family $U\cap\sigma$ has size below $\eta$. Its intersection belongs
to the proper ultrafilter $U$, so is nonempty. Let
\[
b=\min\bigcap(U\cap\sigma).
\]
For every $A\in\sigma$ with $A\subseteq\tau$, elementarity puts
$\tau\setminus A$ in $\sigma$. The ultrafilter alternatives give
\begin{equation}\label{eq:trace}
A\in U\quad\Longleftrightarrow\quad b\in A.
\end{equation}
Now $U$ is the unique ultrafilter $W\in\sigma$ on $\tau$ with this
trace. If two distinct candidates in $\sigma$ had the same trace,
elementarity would supply a set $A\in\sigma$ on which they differ,
contradicting \eqref{eq:trace}. The resulting definition uses only
$\sigma$ and the ordinals $b,\tau$ (and, if satisfaction is used,
$\alpha$). It therefore makes $U$ ordinal definable from $\sigma$.
Neither $b\in\sigma$ nor principality of $U$ is required; uniqueness
is only among the candidates inside $\sigma$.
\end{proof}

\begin{proposition}[Inherited barrier, with proof]\label{prop:barrier}
Assume $\CEx_\gamma(\lambda)$. For every cardinal $\eta\ge q(\gamma)$,
\[
\cf^V(\lambda)=\omega,
\qquad
\cf^{\HCD(\eta)}(\lambda)=\lambda.
\]
Moreover no short cofinal subset of $\lambda$ belongs to $\CD(\eta)$.
\end{proposition}
\begin{proof}
First let $\eta>\gamma$ and let $x\in\CD(\eta)$, witnessed by finitely
many $\eta$-complete ultrafilters $U_0,\ldots,U_{m-1}$ on ordinals.
Input~\ref{in:BG} supplies an elementary hull $\sigma\prec V_\alpha$
containing these ultrafilters and their underlying ordinals, with
$|\sigma|\le\gamma$ and $\REx(\lambda,\sigma)$. This follows by
intersecting the stationary set there with the club of elementary
hulls containing the listed parameters.

For each $i$, Lemma~\ref{lem:trace} applies because
$|\sigma|\le\gamma<\eta$. Substitution in the definition of $x$ now
gives $x\in\OD_{\{\sigma\}}$. Definability transfer in
Input~\ref{in:BG} gives $\REx(\lambda,x)$. Lemma~\ref{lem:cofinal}
therefore excludes every cofinal $x\subseteq\lambda$ of order type
below $\lambda$.

Put $N=\HCD(\eta)$. The cardinal $\lambda$ is still a cardinal in $N$.
If $N$ contained a cofinal map $f:\mu\to\lambda$ with $\mu<\lambda$,
its range $a\in N\subseteq\CD(\eta)$ would have ambient size below
$\lambda$. As a subset of the initial ordinal $\lambda$, it would
have order type below $\lambda$, a contradiction. This reasoning needs
only $\ZF$ in $N$. Ambient countable cofinality follows from
Lemma~\ref{lem:critical}.

Finally suppose $\gamma$ is singular. Every $\gamma$-complete filter
is $\gamma^+$-complete: a family of size $\gamma$ can be partitioned
into $\cf(\gamma)$ many subfamilies of size below $\gamma$; intersect
within each subfamily and then intersect the resulting
$\cf(\gamma)<\gamma$ members. Thus
$\CD(\gamma)=\CD(\gamma^+)$ and
$\HCD(\gamma)=\HCD(\gamma^+)$. This includes the endpoint
$\eta=q(\gamma)=\gamma$; monotonicity gives all larger $\eta$.
\end{proof}

\section{The missing amenability observation}\label{sec:amenability}

For an ordinal $\alpha$ of uncountable cofinality define
\[
E^\alpha_\omega=\{\xi<\alpha:\cf^V(\xi)=\omega\},
\qquad
\calF^V_{\alpha,\omega}
=\{A\subseteq\alpha:\exists C\text{ club in }\alpha\
(C\cap E^\alpha_\omega\subseteq A)\}.
\]
This is the filter generated by the ambient club filter and
$E^\alpha_\omega$. A set $A\subseteq\alpha$ is positive for this
filter precisely when $A\cap E^\alpha_\omega$ is stationary in
$\alpha$. The definition makes sense for singular $\alpha$ as long as
$\cf^V(\alpha)>\omega$.

\begin{definition}[$\omega$-club amenability]\label{def:amenable}
An inner model $N\models\ZFC$ is $\omega$-club amenable in $V$ if,
for every ordinal $\alpha$ with $\cf^V(\alpha)>\omega$,
\[
\calF^V_{\alpha,\omega}\cap N\in N.
\]
This is Goldberg's definition \cite[\S2.2]{dichotomy}. The filter in
this formula is the \emph{ambient} filter, not $N$'s internally
computed club filter.
\end{definition}

The trace expression means the set of members of the filter that
belong to $N$; its members are subsets of $\alpha$. The following
observation applies even before Choice in $N$ is available.

\begin{lemma}[Ordinal-definable subsets of a hereditary class]\label{lem:ODclosure}
If $N=\HCD(\eta)$ and $A$ is an ordinal-definable set with
$A\subseteq N$, then $A\in N$.
\end{lemma}
\begin{proof}
The set $A$ belongs to $\CD(\eta)$ because every ordinal-definable set
does. For each $x\in A$, the assumption $x\in N$ says that every
member of $\tc(\{x\})$ belongs to $\CD(\eta)$. Therefore every member
of $\tc(\{A\})$ belongs to $\CD(\eta)$, including $A$ itself.
This is exactly the condition $A\in\HCD(\eta)$.
\end{proof}

\begin{proposition}[Amenability of the high-completeness models]\label{prop:amenable}
For every infinite cardinal $\eta$ and every ordinal $\alpha$ of
uncountable ambient cofinality,
\[
\calF^V_{\alpha,\omega}\cap\HCD(\eta)\in\HCD(\eta).
\]
In particular, whenever $\HCD(\eta)$ satisfies Choice, it is an
$\omega$-club-amenable inner model of $\ZFC$.
\end{proposition}
\begin{proof}
The trace is a set by Separation from $\calP(\alpha)$. It is definable
in $V$ from the two ordinal parameters $\alpha,\eta$, and all its
members belong to $\HCD(\eta)$. Lemma~\ref{lem:ODclosure} applies.
\end{proof}

\begin{remark}[What this does not say]
The ambient club or stationary sets themselves need not belong to
$N$. Nor does this identify ambient stationarity with internal
stationarity. It says that \emph{for sets already in $N$}, the test
of ambient $\omega$-club positivity can be encoded by a set of $N$.
This is exactly the amount of information required by the seed
construction below.
\end{remark}

\section{Stationary lifting and the local seed lemma}\label{sec:seed}

The method is an explicit version of the stationary-seed arguments in
Goldberg's work \cite[Proposition~2.3 and Lemma~2.7]{dichotomy}; compare
also \cite[Theorem~4.13]{unique} and \cite[Proposition~2.2]{note}.
In contrast to the same-parameter $\HCD$ covering argument, the lower
compactness bound and the higher complete-definability parameter are
separated here. The essential additional bookkeeping is
that the seed belongs to $N^j$, not merely to the ultrapower $M$.
We prove the local form needed here in detail.

\subsection{The continuity calculation}

\begin{lemma}[Lifting a stationary set of countable cofinalities]\label{lem:lift}
Let $\theta$ be an uncountable regular cardinal of $V$ and
$j:V\to M$ a well-founded ultrapower embedding with $j(\omega)=\omega$.
Set $\beta=\sup j``\theta$. If $S\subseteq E^\theta_\omega$ is
stationary in $V$, then $j(S)\cap\beta$ meets every club in $\beta$
that belongs to $M$.
\end{lemma}
\begin{proof}
Fix such a club $C$. Closedness and unboundedness of $C$ in the ordinal
$\beta$ are absolute to $V$. Define, externally in $V$,
\[
f(\xi)=\sup_{\zeta<\xi}j(\zeta)\qquad(\xi<\theta).
\]
At limit indices this function is continuous. Let
\[
D=\{\xi<\theta:\xi\text{ is a nonzero limit ordinal and }f(\xi)\in C\}.
\]
We verify directly that $D$ is club in $\theta$.

For unboundedness, start above any chosen $\xi_0<\theta$ and recursively
choose increasing $\xi_n<\theta$ and $c_n\in C$ so that
\[
j(\xi_n)<c_n<j(\xi_{n+1}).
\]
The choices are possible because both $C$ and $j``\theta$ are
unbounded in $\beta$. Regularity of $\theta$ gives
$\xi=\sup_n\xi_n<\theta$. The two interleaved sequences have the same
supremum, namely $f(\xi)=\sup_n c_n<\beta$; closedness of $C$ puts
$f(\xi)$ in $C$. Thus $\xi\in D$ and $\xi>\xi_0$.
For closedness, if $D\cap\xi$ is unbounded in a limit
$\xi<\theta$, continuity gives
$f(\xi)=\sup\{f(\zeta):\zeta\in D\cap\xi\}\in C$.
The values here are unbounded in $f(\xi)$, so closedness applies.

Choose $\xi\in S\cap D$. A countable increasing cofinal sequence in
$\xi$ is carried by $j$ to a countable cofinal sequence in $j(\xi)$,
because $j(\omega)=\omega$. Consequently
$j(\xi)=\sup j``\xi=f(\xi)\in C$. Also
$j(\xi)\in j(S)$ and $j(\xi)<\beta$. This proves the assertion.
\end{proof}

\subsection{Constructing a seed inside the inner model}

\begin{theorem}[Local stationary-seed covering]\label{thm:localseed}
Let $N\models\ZFC$ be a set-parameter-definable, $\omega$-club-amenable
inner model. Suppose $\delta$ is strongly compact in $V$,
$\mu$ is an ordinal, and $\theta>\max\{\delta,\mu\}$ is regular in $V$.
Assume there is a sequence
\[
\vec A=\langle A_\xi:\xi<\mu\rangle\in N
\]
of pairwise disjoint subsets of $\theta$, each positive for
$\calF^V_{\theta,\omega}$. Then there is a $\delta$-complete fine
ultrafilter $U$ on $P_\delta(\mu)$ such that
\begin{equation}\label{eq:concentrates}
N\cap P_\delta(\mu)\in U.
\end{equation}
In particular every $x\subseteq\mu$ of ambient size below $\delta$
has a $\delta$-cover in $N$.
\end{theorem}
\begin{proof}
Choose a $\delta$-complete fine ultrafilter $U_0$ on
$P_\delta(\theta)$ and form its well-founded ultrapower $j:V\to M$.
Let $u=[\id]_{U_0}$ be the canonical seed. Then
\begin{equation}\label{eq:bigseed}
j``\theta\subseteq u\subseteq j(\theta),
\qquad |u|^M<j(\delta),
\qquad \crit(j)\ge\delta.
\end{equation}
Put $\beta=\sup j``\theta$. Since $j(\theta)$ is regular in $M$ and
$u$ is $M$-small, $\beta<j(\theta)$.

\textbf{Step 1: the cofinality bounds.}
The increasing sequence $j``\theta$ shows $\cf^V(\beta)=\theta$.
To see equality rather than just an upper bound, any cofinal subset of
$\beta$ of size below the regular cardinal $\theta$ could be bounded
by fewer than $\theta$ many corresponding points of $j``\theta$,
contradicting cofinality of that sequence. Hence
$\cf^M(\beta)\ge\theta>\omega$.
On the other hand $u\cap\beta\in M$ is cofinal in $\beta$ and has
$M$-size below $j(\delta)$. Thus
\begin{equation}\label{eq:betacf}
\omega<\cf^M(\beta)<j(\delta).
\end{equation}
Fix in $M$ a club $C\subseteq\beta$ of size
$\cf^M(\beta)<j(\delta)$.

\textbf{Step 2: the seed is definable in $N^j$.}
Let $\vec B=j(\vec A)=\langle B_\xi:\xi<j(\mu)\rangle$.
It belongs to $N^j$ and its entries are pairwise disjoint. By the
transferred amenability statement and \eqref{eq:betacf},
\[
F=\calF^M_{\beta,\omega}\cap N^j\in N^j.
\]
For each $\xi<j(\mu)$, both $B_\xi\cap\beta$ and its complement
in $\beta$ belong to $N^j$. Therefore Separation \emph{inside $N^j$}
constructs the set
\begin{equation}\label{eq:smallseed}
s=\{\xi<j(\mu):\beta\setminus(B_\xi\cap\beta)\notin F\}\in N^j.
\end{equation}
Equivalently, $\xi\in s$ exactly when $B_\xi\cap\beta$ is positive
for the ambient-to-$M$ filter $\calF^M_{\beta,\omega}$.
This is the point where the trace belongs to the smaller model; mere
external definability of $s$ in $M$ would not suffice.

\textbf{Step 3: fineness and smallness of the seed.}
Fix $\xi<\mu$. Since $A_\xi$ is positive,
$S=A_\xi\cap E^\theta_\omega$ is stationary in $V$.
Lemma~\ref{lem:lift} shows that $j(S)\cap\beta$ meets every club of
$\beta$ in $M$. Every member of this set has $M$-cofinality $\omega$,
since $j(S)\subseteq(E^{j(\theta)}_\omega)^M$. It is contained in
$B_{j(\xi)}\cap\beta$. Therefore $j(\xi)\in s$, and we have
\begin{equation}\label{eq:seedfine}
j``\mu\subseteq s.
\end{equation}
Every $\xi\in s$ satisfies $B_\xi\cap C\ne\varnothing$. Choosing
its least member defines in $M$ an injection
\[
\xi\longmapsto\min(B_\xi\cap C)
\quad\text{from }s\text{ into }C,
\]
because the $B_\xi$ are disjoint. Hence $|s|^M<j(\delta)$.
Together with $s\subseteq j(\mu)$, this says
$s\in j(P_\delta(\mu))$.

\textbf{Step 4: derive the ultrafilter.}
Define in $V$
\[
U=\{Z\subseteq P_\delta(\mu):s\in j(Z)\}.
\]
Since $s\in j(P_\delta(\mu))$, elementarity proves the ultrafilter
axioms. The critical point bound in \eqref{eq:bigseed} proves
$\delta$-completeness: for a sequence of fewer than $\delta$ many
members, $j$ fixes the index ordinal and the seed lies in every entry
of its image, hence in the image of the intersection.
Equation~\eqref{eq:seedfine} proves fineness. Finally the definition of
$N^j$ gives
\[
j(N\cap P_\delta(\mu))=N^j\cap j(P_\delta(\mu)).
\]
Our seed belongs to the right-hand side by \eqref{eq:smallseed}, proving
\eqref{eq:concentrates}.

\textbf{Step 5: extract a cover.}
For $x\subseteq\mu$ with $|x|<\delta$, intersect the concentrating
set \eqref{eq:concentrates} with the fine cones for all $\xi\in x$.
There are fewer than $\delta$ cones, so the intersection is in the
proper ultrafilter $U$ and is nonempty. A member $a$ belongs to $N$,
contains $x$, and has ambient size below $\delta$.
Lemma~\ref{lem:smallness} converts this to $|a|^N<\delta$.
\end{proof}

\begin{remark}[No unnecessary closure of the ultrapower]
The proof does not assume $M^\theta\subseteq M$ or
$j``\theta\in M$. The small set $u\in M$ covers that image;
$\beta$ is an ordinal and hence belongs to $M$.
The lower bound on its cofinality is computed externally, while the
upper bound is witnessed by $u\cap\beta\in M$. This is why strong
compactness suffices where supercompactness might otherwise be
mistakenly assumed.
\end{remark}

\section{Amenable grounds cannot regularize these ordinals}\label{sec:grounds}

\subsection{Why a ground supplies sufficiently high stationary families}

\begin{lemma}[High stationary partitions in a ground]\label{lem:groundpartition}
Suppose $V=W[G]$ for a set forcing $\PP\in W$, and $\mu$ is any
ordinal. There is an ordinal $\theta>\mu$ that is a successor cardinal
in $W$ and regular in $V$, together with a sequence
$\langle S_\xi:\xi<\mu\rangle\in W$ of pairwise disjoint subsets of
$(E^\theta_\omega)^W$, each stationary in $V$. The ordinal $\theta$
can be chosen above any additional prescribed bound.
\end{lemma}
\begin{proof}
In $W$, choose an infinite cardinal $\nu$ greater than $\mu$ and
$|\PP|^W$, and beyond the prescribed bound. Put $\theta=(\nu^+)^W$.
Then $W$ regards $\PP$ as $\theta$-cc. Such forcing preserves
regularity of $\theta$ and all $W$-stationary subsets of $\theta$;
a direct proof is given in Appendix~\ref{app:forcing}.

Inside $W$, the set $(E^\theta_\omega)^W$ is stationary. The Ulam-matrix
argument at a successor cardinal, proved in
Appendix~\ref{app:splitting}, splits it into $\theta$ many disjoint
stationary subsets. Take the first $\mu$ cells of the resulting
sequence. They remain stationary in $V$. A countable cofinal sequence
witnessing membership in $(E^\theta_\omega)^W$ remains a countable
cofinal sequence in $V$, so all these cells are subsets of the actual
$E^\theta_\omega$.
\end{proof}

\begin{theorem}[Covering by an $\omega$-club-amenable ground]\label{thm:groundcover}
Let $W\models\ZFC$ be a set-parameter-definable ground of $V$ that is
$\omega$-club amenable in $V$. If $\delta$ is strongly compact in $V$,
then $W$ has the $\delta$-cover property in $V$.
\end{theorem}
\begin{proof}
First let $x\subseteq\mu$ have size below $\delta$. Apply
Lemma~\ref{lem:groundpartition}, choosing the regular $\theta$ also
above $\delta$. The resulting sequence is a family of positive sets
for $\calF^V_{\theta,\omega}$, so
Theorem~\ref{thm:localseed} gives a cover $a\in W\cap P_\delta(\mu)$
of $x$. Lemma~\ref{lem:smallness} gives $|a|^W<\delta$.

For the usual cover property on arbitrary sets $x\subseteq W$, choose
an ordinal $\alpha$ with $x\subseteq V_\alpha^W$. Choice in $W$
provides a bijection $b:\mu\to V_\alpha^W$ for some ordinal $\mu$.
The inverse image $b^{-1}[x]$ is an ambient set of ordinals of size
below $\delta$. Cover it by $a\in W$ of $W$-size below $\delta$ as
just proved. Then $b[a]\in W$ covers $x$ and still has $W$-size below
$\delta$.
\end{proof}

\begin{corollary}[No amenable-ground regularization]\label{cor:noregularization}
Suppose $\delta$ is strongly compact in $V$, $\lambda>\delta$ is an
ordinal with $\cf^V(\lambda)=\omega$, and $W$ is a set-parameter-definable,
$\omega$-club-amenable inner model of $\ZFC$ in which $\lambda$ is
regular. Then $W$ is not a set-forcing ground of $V$.
\end{corollary}
\begin{proof}
If it were a ground, take a countable cofinal $c\subseteq\lambda$ in
$V$. Theorem~\ref{thm:groundcover} would give $a\in W$ containing $c$
with $|a|^W<\delta$. Replace $a$ by $a\cap\lambda$. This set is
cofinal in $\lambda$ because it contains $c$. Cofinality of a fixed set
of ordinals is absolute between transitive models with the same
ordinals. Thus $W$ contains a cofinal subset of its regular cardinal
$\lambda$ of $W$-size below $\delta<\lambda$, a contradiction.
\end{proof}

This argument is independent of the size, distributivity, or chain
condition of the forcing at $\lambda$ itself. Every \emph{set} forcing
has a high enough level at which it preserves the stationary partition
needed by the proof. This is what makes the result stronger than a
bound on a singularizing forcing at one specified successor.

\section{The conditional inconsistency theorem}\label{sec:mainproof}

\begin{theorem}[No high-completeness ground]\label{thm:noHCDground}
Assume $\delta<\lambda\le\gamma$, $\SC(\delta)$, and
$\CEx_\gamma(\lambda)$. Then for every cardinal $\eta\ge q(\gamma)$,
$\HCD(\eta)$ is not a set-forcing ground of $V$.
\end{theorem}
\begin{proof}
Suppose $N=\HCD(\eta)$ were a ground. In particular it would satisfy
$\ZFC$, whether or not Choice was known beforehand.
Proposition~\ref{prop:amenable} makes it $\omega$-club amenable, and it
is definable from the ordinal $\eta$. By
Proposition~\ref{prop:barrier}, $\lambda$ is regular in $N$ while
$\cf^V(\lambda)=\omega$. Since $\delta<\lambda$ is strongly compact
in $V$, Corollary~\ref{cor:noregularization} gives the contradiction.
\end{proof}

\begin{proof}[Proof of Theorem~\ref{thm:main}]
Suppose, contrary to the conclusion, that a strongly compact
$\kappa>\gamma$ exists. Then $\kappa\ge q(\gamma)$.
Input~\ref{in:ground} says that $\HCD(\kappa)$ is a set-forcing ground
of $V$. Theorem~\ref{thm:noHCDground} says that it is not.
\end{proof}

\begin{corollary}[A proper class of strongly compact cardinals]\label{cor:properclass}
Assume there is a proper class of strongly compact cardinals, and let
$\delta_0$ be the least one. Every cover-exacting cardinal is strictly
below $\delta_0$. Equivalently, no cardinal at or above $\delta_0$ is
$\gamma$-cover exacting for any cardinal $\gamma$.
\end{corollary}
\begin{proof}
If $\lambda>\delta_0$ were $\gamma$-cover exacting, choose a strongly
compact $\kappa>\gamma$ and apply Theorem~\ref{thm:main} with
$\delta=\delta_0$. Equality $\lambda=\delta_0$ is impossible because
an exacting cardinal has cofinality $\omega$, whereas a strongly
compact cardinal is uncountable regular.
\end{proof}

\begin{corollary}[A bound forced on every cover width]\label{cor:bound}
If a cover-exacting $\lambda$ lies above a strongly compact cardinal,
then the class of strongly compact cardinals is bounded. Every witness
$\gamma$ to $\CEx_\gamma(\lambda)$ is an upper bound for that class.
In particular, a $\lambda$-cover-exacting $\lambda$ cannot have strongly
compact cardinals both below and above it.
\end{corollary}
\begin{proof}
For each fixed cover bound $\gamma$, Theorem~\ref{thm:main} excludes
all strongly compact cardinals greater than $\gamma$. The diagonal
case sets $\gamma=\lambda$.
\end{proof}

\begin{remark}[The strict inequality at the upper endpoint]
The conclusion excludes $\kappa>\gamma$. It does not, by this
argument, exclude $\kappa=\gamma$ when $\gamma$ is regular. At a regular
cover bound the inherited reconstruction requires completeness strictly
greater than $\gamma$. The singular endpoint has already been included
through $q(\gamma)$, but a strongly compact cardinal is not singular.
\end{remark}

\subsection{Cross-check against the published dichotomy}\label{sec:crosscheck}

Goldberg proves the following result
\cite[Lemma~2.7]{dichotomy}: if $N$ is $\omega$-club amenable and
$\delta$ is $\omega_1$-strongly compact, then either all sufficiently
large regular cardinals of $V$ are measurable in $N$, or $N$ has the
$(\omega_1,\delta)$-cover property. The latter means that every
countable subset of $N$ is covered by a set in $N$ of $N$-size below
$\delta$. A strongly compact cardinal satisfies the stated weaker
hypothesis.

Here is a short second verification of the no-ground step. Suppose
$V=N[G]$ by a set forcing. At arbitrarily large successor cardinals of
$N$, the forcing preserves both that successor's cardinality and its
regularity. In particular, there are arbitrarily large regular
cardinals of $V$ that are still successor cardinals in $N$. They cannot
be measurable in $N$, since measurable cardinals are inaccessible in
any model of $\ZFC$. Thus the first alternative is impossible. The
countable-cover alternative follows and contradicts regularity in
$N$ of a $\lambda>\delta$ of ambient cofinality $\omega$, exactly as
in Corollary~\ref{cor:noregularization}.

This verification has the same classical stationary-seed background
as our detailed proof; it is not a claim of a new HOD dichotomy. It is
useful because it checks the main conclusion against an independently
stated published theorem, and uses only countable covering rather than
the full $\delta$-cover property we proved for strongly compact
$\delta$.

\section{Quantitative stationary width and inner measurability}\label{sec:width}

The preceding proof has a local consequence even when no ground is
assumed. It records how extensively high stationary partitions must
fail in a hypothetical remaining model.

\begin{theorem}[A bound on ambient stationary width]\label{thm:width}
Let $N\models\ZFC$ be set-parameter definable and $\omega$-club
amenable. Suppose $\delta$ is strongly compact in $V$,
$\delta<\lambda$, $\cf^V(\lambda)=\omega$, and $\lambda$ is regular
in $N$. For every regular cardinal $\theta>\lambda$ of $V$, the filter
\[
F_\theta=\calF^V_{\theta,\omega}\cap N\in N
\]
is $\lambda$-weakly saturated in $N$: there is no sequence in $N$ of
$\lambda$ pairwise disjoint $F_\theta$-positive subsets of $\theta$.
\end{theorem}
\begin{proof}
Such a sequence would meet all the assumptions of
Theorem~\ref{thm:localseed} with $\mu=\lambda$. A countable cofinal
$c\subseteq\lambda$ would have an $N$-cover of $N$-size below
$\delta$, contradicting regularity of $\lambda$ in $N$.
For a set $A\in N$ with $A\subseteq\theta$, positivity for the filter
object $F_\theta$ as computed in $N$ is exactly ambient
$\calF^V_{\theta,\omega}$-positivity: in both cases it says that
$\theta\setminus A$ is not in the trace. Thus no absoluteness of
stationarity has been assumed in translating the contradiction into
the internal weak-saturation statement.
\end{proof}

\begin{corollary}[How many cells can survive?]\label{cor:survivors}
Under the hypotheses of Theorem~\ref{thm:width}, let
$\vec S=\langle S_\xi:\xi<\tau\rangle\in N$ be a sequence of pairwise
disjoint subsets of $\theta$, where $\theta>\lambda$ is regular in
$V$. Then
\[
I=\{\xi<\tau:S_\xi\text{ is positive for }
\calF^V_{\theta,\omega}\}
\]
belongs to $N$ and has $N$-cardinality below $\lambda$.
In particular, if $\vec S$ is an $N$-partition of
$(E^\theta_\omega)^N$ into $\theta$ stationary cells, fewer than
$\lambda$ of those cells remain stationary in $V$.
\end{corollary}
\begin{proof}
The trace $F_\theta\in N$ and the sequence $\vec S\in N$ define
$I$ by Separation in $N$, using the complement test for positivity.
If $|I|^N\ge\lambda$, Choice in $N$ would select a subsequence of
length $\lambda$, contradicting Theorem~\ref{thm:width}.
For the last assertion, every cell is a subset of the true
$E^\theta_\omega$, so positivity is precisely ambient stationarity.
\end{proof}

For $N=\HCD(\eta)$ and a cover-exacting $\lambda$, the cardinal
$\lambda$ is also an ambient cardinal. In that application
Lemma~\ref{lem:smallness} converts the bound to ambient size below
$\lambda$ as well. Thus a partition of length $\theta>\lambda$ has
$\theta$ many cells that cease to be stationary. This is a restriction
at every applicable high level, not just at $(\lambda^+)^N$.

\begin{lemma}[Internal completeness of the trace]\label{lem:filtercomplete}
If $N$ is $\omega$-club amenable and $\theta$ is uncountable regular in
$V$, then $N$ regards $F_\theta$ as a proper $\theta$-complete filter
containing every co-bounded subset of $\theta$ that belongs to $N$.
\end{lemma}
\begin{proof}
The ordinal $\theta$ remains a regular cardinal in $N$: any smaller
cardinality or cofinality witness in $N$ would remain such a witness
in $V$. The ambient filter is proper because every club in a regular
uncountable $\theta$ meets $E^\theta_\omega$. It contains the
co-bounded sets.

Let $\langle A_i:i<\rho\rangle\in N$ with $\rho<\theta$ and each
$A_i\in F_\theta$. In $V$ choose a club $C_i$ witnessing membership
of each $A_i$ in the ambient filter. Regularity of $\theta$ makes
$C=\bigcap_{i<\rho}C_i$ a club. Then
$C\cap E^\theta_\omega\subseteq\bigcap_{i<\rho}A_i$.
This intersection belongs to $N$ and to the ambient filter, hence to
$F_\theta$. All the filter axioms and the asserted internal
completeness follow. The witness clubs need not belong to $N$.
\end{proof}

\begin{theorem}[An explicit inner-measurability threshold]\label{thm:measurable}
Under the hypotheses of Theorem~\ref{thm:width}, every regular cardinal
$\theta$ of $V$ such that
\[
\theta>(2^\lambda)^N
\]
is measurable in $N$.
\end{theorem}
\begin{proof}
Work inside $N$. The ordinal $\lambda$ is an infinite cardinal,
$F_\theta$ is $\lambda$-weakly saturated by
Theorem~\ref{thm:width}, and it is $\theta$-complete by
Lemma~\ref{lem:filtercomplete}. Since
$\theta>(2^\lambda)^N$, this supplies the
$((2^\lambda)^+)^N$-completeness needed by Ulam's atomicity lemma.
The lemma, proved in Appendix~\ref{app:atomicity}, partitions
$\theta$ modulo the dual ideal into fewer than $\lambda$ positive
atoms. Restricting the filter to any one of these atoms yields a
$\theta$-complete ultrafilter in $N$ on $\theta$. It is nonprincipal:
the filter contains the complement of every singleton. Thus $N$
regards $\theta$ as measurable.
\end{proof}

\begin{corollary}[The high-$\HCD$ conclusion without an upper compact]\label{cor:HCDmeasurable}
Assume $\SC(\delta)$, $\delta<\lambda\le\gamma$, and
$\CEx_\gamma(\lambda)$. Let $\eta\ge q(\gamma)$ and assume explicitly
that $N=\HCD(\eta)$ satisfies $\ZFC$. Then
Theorem~\ref{thm:width} and Corollary~\ref{cor:survivors} hold for $N$,
and every ambient regular $\theta>(2^\lambda)^N$ is measurable in $N$.
In particular this holds for every ambient regular
$\theta\ge(2^\lambda)^{+V}$.
\end{corollary}
\begin{proof}
Propositions~\ref{prop:barrier} and~\ref{prop:amenable} supply the
hypotheses of Theorem~\ref{thm:width}. Apply
Theorem~\ref{thm:measurable}.
For the final bound, put $\chi=(2^\lambda)^N$. A bijection in $N$
from $\chi$ to $(\calP(\lambda))^N$ remains an injection in $V$ into
$\calP^V(\lambda)$. Hence $|\chi|^V\le 2^\lambda$ and
$\chi<(2^\lambda)^{+V}$. This is an ordinal comparison; it does not
assert that the two models calculate $2^\lambda$ identically.
\end{proof}

\begin{remark}[A second route from the quantitative theorem to no grounds]
A set-forcing ground has arbitrarily high successor cardinals that
remain regular in the extension. Choose one above $(2^\lambda)^N$.
Theorem~\ref{thm:measurable} would make it measurable in $N$, even
though it is a successor there. This is another way to see the ground
obstruction. The direct covering proof is shorter and does not need
filter atomicity; the atomicity argument explains the stronger
structure forced on any remaining high inner model.
\end{remark}

\section{A weaker lower compactness assumption}\label{sec:weaker}

The contradiction only needs to cover a countable cofinal set.
Consequently one may weaken the lower strongly compact hypothesis
without changing the upper one.

\begin{definition}[The countably complete fine-ultrafilter hypothesis]
For an uncountable cardinal $\delta$, let $\Ffine(\delta)$ mean:
for every sufficiently large regular $\theta$ there is a countably
complete fine ultrafilter on $P_\delta(\theta)$.
\end{definition}

\begin{proposition}[The countable version of the seed argument]\label{prop:weakseed}
Replace strong compactness in Theorem~\ref{thm:localseed} by the
existence of a countably complete fine ultrafilter on
$P_\delta(\theta)$. Its conclusion then holds with ``countably
complete'' in place of ``$\delta$-complete,'' and every countable
$x\subseteq\mu$ has a cover in $N$ of $N$-size below $\delta$.
\end{proposition}
\begin{proof}
The well-founded ultrapower still has the seed $u$ with
$j``\theta\subseteq u$ and $|u|^M<j(\delta)$. The cofinality bounds,
the transferred filter trace, the set $s\in N^j$, its $M$-smallness,
and inclusion $j``\mu\subseteq s$ are proved exactly as in
Theorem~\ref{thm:localseed}. They do not use
$\crit(j)\ge\delta$. Continuity at countable cofinality uses only
$j(\omega)=\omega$, which holds for the well-founded ultrapower.

The derived ultrafilter is countably complete: given a countable
sequence of its members, $j$ fixes the index set $\omega$, so $s$
lies in the image of their intersection. For a countable $x$, only
countably many fine cones must be intersected with the concentrating
set. This produces the required cover; Lemma~\ref{lem:smallness}
again supplies the internal size bound. No claim of
$\delta$-completeness is made in this version.
\end{proof}

\begin{theorem}[Weakened no-sandwich theorem]\label{thm:weakmain}
There are no cardinals $\delta<\lambda\le\gamma<\kappa$ such that
$\Ffine(\delta)$, $\CEx_\gamma(\lambda)$, and $\SC(\kappa)$ all hold.
More generally, under $\Ffine(\delta)$ and
$\CEx_\gamma(\lambda)$ with $\delta<\lambda$, no
$\HCD(\eta)$ for $\eta\ge q(\gamma)$ is a set-forcing ground of $V$.
\end{theorem}
\begin{proof}
If such a high model were a ground, choose $\theta$ in
Lemma~\ref{lem:groundpartition} large enough for $\Ffine(\delta)$
to apply. Proposition~\ref{prop:weakseed} covers a countable cofinal
subset of $\lambda$ by an internally $<\delta$ set in that model.
The inherited regularity barrier contradicts this. An upper strongly
compact $\kappa$ then supplies the forbidden ground by
Input~\ref{in:ground}.
\end{proof}

For a familiar sufficient hypothesis, let $\delta$ be regular and
$\omega_1$-strongly compact, in the sense that every $\delta$-complete
filter extends to a countably complete ultrafilter. On
$P_\delta(\theta)$, the filter generated by the sets
\[
\{b\in P_\delta(\theta):a\subseteq b\}
\qquad(a\in P_\delta(\theta))
\]
is proper and $\delta$-complete: regularity of $\delta$ ensures that
the union of fewer than $\delta$ many such $a$'s still has size below
$\delta$. A countably complete ultrafilter extending it is fine.
Thus this regular $\omega_1$-strong compactness assumption implies
$\Ffine(\delta)$ and can replace the lower strong compactness in the
main result. We do not assert an equivalence between these
formulations at singular $\delta$.

\section{Scope, proof audit, and remaining questions}\label{sec:audit}

\subsection{What has and has not been established}

The main theorem is a conditional inconsistency of a specific
\emph{joint} large-cardinal configuration. It is not merely an
inconsistency obtained by appending the Ground Axiom or an unproved
stabilization principle: both were removed. The upper and lower
strongly compact cardinals supply different, indispensable inputs to
the proof. The upper one gives a high-completeness ground; the lower
one supplies a fine ultrafilter used to derive countable covering in
that ground.

The bare theory
\[
\ZFC+\exists\delta<\lambda\le\gamma\,
\bigl(\SC(\delta)\wedge\CEx_\gamma(\lambda)\bigr)
\]
is not refuted by this report. Problem~5.2 in the Blue--Goldberg notes
asks about precisely that remaining configuration \cite{cover}.
Our conclusion is that any such model must have no strongly compact
above its cover bound, no high-$\HCD$ ground, and---whenever a high
$\HCD$ model satisfies Choice---the stationary-width and measurability
behavior of Section~\ref{sec:width}.

Nor does the argument apply automatically to ordinary exactingness,
to ultraexactingness without a cover hypothesis, or to a hypothesis
whose cover bound changes with the height. The fixed bound $\gamma$
is used in the trace reconstruction and is the reason an upper
$\kappa>\gamma$ can be selected uniformly.

\subsection{Dependency audit}

\begin{center}
\small
\renewcommand{\arraystretch}{1.22}
\begin{tabularx}{\textwidth}{@{}P{.26\textwidth}YY@{}}
\toprule
\textbf{Ingredient}&\textbf{Source or proof}&\textbf{Role}\nobreak\\\midrule
Completeness barrier&Inherited from the supplied synthesis; reconstructed in Section~\ref{sec:barrier} from Inputs A--B.&Makes $\lambda$ regular in the high model.\\
High $\HCD$ is a ground&Goldberg, Input~\ref{in:ground}; requires the upper strongly compact.&Converts the inner model into a set-forcing ground.\\
$\omega$-club amenability&Proposition~\ref{prop:amenable}, by hereditary definability.&Places the stationary-positivity trace in the high model.\\
Seed belongs to $N^j$&Equation~\eqref{eq:smallseed}; Separation from the trace.&Makes the derived measure concentrate on $N$.\\
Seed is small&Equations~\eqref{eq:betacf}--\eqref{eq:seedfine}; disjointness gives an injection into a small club.&Provides $<\delta$ covers.\\
High stationary family&Lemma~\ref{lem:groundpartition} and Appendix A.&Uses only that the forcing is a set.\\
No-sandwich contradiction&Theorems~\ref{thm:noHCDground} and~\ref{thm:main}.&New application; no stabilization or Ground Axiom.\\
Quantitative measurability&Theorem~\ref{thm:width}, trace completeness, and Appendix B.&A stronger local consequence, not needed for the main contradiction.\\
\bottomrule
\end{tabularx}
\end{center}

\subsection{Sensitive points checked explicitly}

The inequality in trace reconstruction is
$|\sigma|\le\gamma<\eta$; it is not $|\sigma|<\gamma$.
The ultrafilters used to define high-$\HCD$ sets have completeness
above the cover bound, whereas the stationary-seed ultrafilter comes
from the \emph{lower} compact cardinal and has completeness
$\delta$. They are different ultrafilters with different jobs.

The model $N^j$ is the interpretation in $M$ of the definition of $N$.
It is not $j``N$. The club-filter trace used there is calculated in
$M$, rather than in $N^j$. Its membership in $N^j$ follows by
transferring the actual amenability statement. Merely saying that
stationarity is definable would not justify the proof. The ordinal
$\beta$ need not be regular in $M$; amenability was required at every
ordinal of uncountable cofinality, not just at regular cardinals.

No ground stationarity is assumed to be preserved at a singularized
ordinal or its first old successor. The proof moves above the size of
the forcing. The Ulam-matrix and chain-condition proofs in Appendix A
show exactly why that move is legitimate.

The upper compact cardinal is assumed strongly compact in $V$, not
merely in a ground. The lower compact cardinal is likewise an ambient
hypothesis. Choice in a high $\HCD$ model is obtained from its being a
ground or is explicitly assumed; it is not inherited from the ambient
Axiom of Choice alone.

\subsection{Research questions left by the result}

The central remaining question is whether the lower-compact,
cover-exacting configuration can exist when the strongly compact
cardinals are bounded. A more focused obstruction would be to derive
that some $\HCD(\eta)$ with $\eta\ge q(\gamma)$ is a ground without
assuming an upper strongly compact. The present work proves that any
such derivation would finish the inconsistency argument, but does not
supply it.

Another precise target is the Choice question for the high-completeness
models under the lower-compact hypothesis alone. The measurable-cardinal
conclusion above is conditional on Choice there. Eliminating that
qualification would require an argument about the inner model, not a
change of notation in the stationary-seed proof.

The supplied Lean sources were used as a reference for the hypotheses
and the inherited barrier. No new Lean theorem is included, and no
compilation of those sources is represented as a verification of the
new result. The present proof obligations are conventional mathematical
ones, isolated in the table above. The bibliographic search and source
checks support the attribution of the inputs, not a claim that no one
has previously noticed this consequence.

\appendix
\section{Classical stationary facts used in the ground argument}\label{app:stationary}

\subsection{Splitting at a successor by an Ulam matrix}\label{app:splitting}

\begin{lemma}[Successor splitting]\label{lem:ulam}
In $\ZFC$, if $\nu$ is infinite, $\theta=\nu^+$, and
$S\subseteq\theta$ is stationary, then $S$ has a family of $\theta$
pairwise disjoint stationary subsets. The family can be made a
partition of $S$.
\end{lemma}
\begin{proof}
For every $\beta\in[\nu,\theta)$ choose a bijection
$e_\beta:\nu\to\beta$. For $\xi<\theta$ and $i<\nu$, put
\[
A_i^\xi=\{\beta<\theta:\beta>\max\{\nu,\xi\}
\text{ and }e_\beta(i)=\xi\}.
\]
For fixed $\xi$, the sets $A_i^\xi$ as $i$ varies partition a tail of
$\theta$. At least one $S\cap A_i^\xi$ is stationary: otherwise
intersect the fewer than $\theta$ many clubs witnessing
nonstationarity, contradicting stationarity of $S$. Choose one such
$i=i(\xi)$.

There is a fixed $i_*<\nu$ for which
$B=\{\xi<\theta:i(\xi)=i_*\}$ has size $\theta$. If every fiber had
size below the regular cardinal $\theta$, their union over
$\nu<\theta$ many indices would have size below $\theta$, although
that union is $\theta$.
For distinct $\xi,\zeta\in B$, the sets
$A_{i_*}^\xi$ and $A_{i_*}^\zeta$ are disjoint, because
$e_\beta(i_*)$ cannot have two values. Enumerating $B$ gives the
required stationary family. Any unused part of $S$ can be added to
one cell to obtain a partition, without affecting disjointness or
stationarity.
\end{proof}

The proof is applied \emph{inside the ground}. At a regular
uncountable $\theta$ there, $E^\theta_\omega$ is stationary: given a
club, build a strictly increasing countable sequence from it and take
its supremum, which is still below $\theta$ and belongs to the club.
The supremum has cofinality $\omega$.

\subsection{Chain-condition preservation, with names}\label{app:forcing}

\begin{lemma}[$\theta$-cc preservation]\label{lem:ccpreserve}
Suppose $W\models\ZFC$, $\theta$ is regular uncountable in $W$, and
$\PP\in W$ is $\theta$-cc in $W$. Then forcing with $\PP$ preserves
regularity of $\theta$ and every stationary $S\subseteq\theta$ of
$W$. It also preserves all $W$-cardinals at or above $\theta$.
\end{lemma}
\begin{proof}
All choices of antichains in this proof are made in $W$.
For a name $\dot f:\rho\to\theta$ with $\rho<\theta$, a maximal
antichain deciding each $\dot f(i)$ has size below $\theta$.
Regularity of $\theta$ bounds the union of all possible values for
all $i<\rho$ below $\theta$. Thus no new such map can be cofinal,
and $\theta$ stays regular.

Now let $p$ force that $\dot C$ is club in $\theta$. Work below $p$.
For each $\alpha<\theta$, choose a maximal antichain $B_\alpha$
whose members decide an ordinal $\beta_q>\alpha$ belonging to
$\dot C$. This is possible because $\dot C$ is forced unbounded.
The chain condition gives
\[
f(\alpha)=\sup\{\beta_q+1:q\in B_\alpha\}<\theta.
\]
Let $D\subseteq\theta$ be the ground-model club of nonzero limit
ordinals $\xi$ closed under $f$, meaning
$f(\alpha)<\xi$ for all $\alpha<\xi$.
In any generic extension containing $p$, for each $\alpha<\xi\in D$
the generic meets $B_\alpha$, so $\dot C$ has a point strictly between
$\alpha$ and $\xi$. Closedness therefore gives $\xi\in\dot C$.
Thus $p$ forces $D\subseteq\dot C$.
If $S$ is stationary in $W$, choose $\xi\in S\cap D$ in $W$.
Then $p$ forces that $S$ meets $\dot C$. Since the name and condition
were arbitrary, $S$ remains stationary.

For cardinal preservation, suppose $\chi\ge\theta$ is a $W$-cardinal
and $\rho<\chi$. In a name for a function from $\rho$ to $\chi$,
there are fewer than $\theta$ possible values per argument, by maximal
antichains. If $\rho<\theta$, their union has size below $\theta$;
if $\rho\ge\theta$, it has $W$-size at most $|\rho|^W\cdot\theta
$ and hence below $\chi$. In neither case can the function be onto
$\chi$. No cardinal at or above $\theta$ is collapsed.
\end{proof}

\section{The filter-atomicity lemma, with proof}\label{app:atomicity}

This is the classical Ulam lemma used in
\cite[Lemma~2.1]{dichotomy}. It is included so that the quantitative
measurability consequence does not rest on an unexplained atomicity
step. The lemma is applied inside $N$, where Choice is available.

\begin{lemma}[Ulam atomicity]\label{lem:atomicity}
Let $\lambda$ be an infinite cardinal and $F$ a proper filter on a set
$X$. Suppose $F$ is $(2^\lambda)^+$-complete and there is no family of
$\lambda$ pairwise disjoint $F$-positive sets. Then $X$ can be
partitioned modulo the dual ideal into fewer than $\lambda$ positive
atoms. Consequently $F$ is an intersection of fewer than $\lambda$
ultrafilters, each inheriting every degree of completeness of $F$.
\end{lemma}
\begin{proof}
Let $I=\{A\subseteq X:X\setminus A\in F\}$ be the dual ideal and
put $\rho=(2^\lambda)^+$. A positive set $A$ is an \emph{atom} when
for each $B\subseteq A$, either $B\in I$ or $A\setminus B\in I$.
A positive non-atom can therefore be split into two disjoint positive
sets.

Suppose no partition modulo $I$ into fewer than $\lambda$ positive
atoms exists. We construct refining partitions modulo $I$ through
stages $\alpha\le\lambda$. Start with $\{X\}$. At a successor stage,
choose a non-atom cell and split it into two positive cells, leaving
the others alone. There must be a non-atom cell by the supposition.
Every partition we obtain has fewer than $\lambda$ positive cells,
by the assumed weak saturation.

At a limit stage $\alpha\le\lambda$, take all intersections along
coherent paths through the preceding partitions. There are at most
$\lambda^\lambda=2^\lambda$ such paths. Their intersections partition
what has not already been discarded as ideal-small. Discard every
intersection that is in $I$. The union of all discarded sets remains
in $I$: there are at most $2^\lambda<\rho$ candidates at each stage
and at most $\lambda$ stages, and $I$ is $\rho$-complete. Thus the
positive intersections still form a partition modulo $I$, again with
fewer than $\lambda$ cells. They refine the preceding partitions by
actual inclusion. This justifies the recursion, including its final
stage $\lambda$.

Let $\calD$ be that final partition; it has size $\tau<\lambda$.
Every positive cell at every earlier stage contains a positive member
of $\calD$. Indeed $\calD$ covers $X$ modulo $I$ and refines that
stage; a positive cell with no final positive descendant would be
contained in the discarded ideal-small set.
For each of the $\lambda$ successor-stage splits, choose one final
cell inside each of its two children. This associates to the split an
unordered pair of distinct members of $\calD$.
Pairs associated to different splits must differ: the two members of
such a pair are in the same cell immediately before that split and
in distinct cells afterward, so the split is the unique first stage
at which they are separated. We have obtained an injection from
$\lambda$ into $[\calD]^2$, impossible because
$|[\calD]^2|\le\tau^2<\lambda$. The recursion must therefore have
reached a partition into positive atoms before this contradiction.

Let $\{A_i:i<\tau\}$, $\tau<\lambda$, be such a partition modulo
$I$. For each $i$ define
\[
U_i=\{B\subseteq X:A_i\setminus B\in I\}.
\]
The atom property gives the ultrafilter alternatives, and $U_i$ is
proper since $A_i\notin I$. If $I$ is closed under unions of fewer
than a cardinal $\zeta$, then $U_i$ is $\zeta$-complete: intersecting
fewer than $\zeta$ members has complement inside $A_i$ contained in
the corresponding union of $I$-small complements.
Finally $F=\bigcap_{i<\tau}U_i$. One inclusion is immediate. For the
other, if all $A_i\setminus B$ are in $I$, then
\[
X\setminus B\subseteq
\left(X\setminus\bigcup_{i<\tau}A_i\right)
\cup\bigcup_{i<\tau}(A_i\setminus B)\in I,
\]
using $\tau<\lambda<\rho$. Hence $B\in F$, as required.
\end{proof}

\section{Source map and reproducibility}\label{app:sources}

The starting archive contains the two reports
\nolinkurl{docs/Large_Cardinals_Synthesis.tex} and
\nolinkurl{docs/Large_Cardinals_Unified_Report.tex}, together with the
\nolinkurl{Cardinals/} Lean development. The specific inherited
material used here is the synthesis's completeness barrier, its
relative-predicate lemmas, and its high-$\HCD$ ground discussion.
The new argument does not depend on the archive's more speculative
Prikry consistency refinement or on its ultraexacting uniformization
results.

Primary sources were checked on 18 September 2026. For the ground
theorem, the exact version cited is arXiv:2103.13961v2. For the
$\omega$-club-amenable covering dichotomy, the checked author PDF is
dated 10 August 2021 and has the result as Lemma~2.7. Its fuller
singular-cardinal dichotomy is Corollary~2.9; only Lemma~2.7 is needed
for our short cross-check. For cover exactingness, the checked notes
are dated 1 July 2026. These version details matter because some
statement numbers, and the strength of hypotheses used in short
summaries, differ.

The accompanying source is self-contained as a \LaTeX{} document:
its bibliography is embedded and it needs no separate figure or
bibliography files. Compile with \texttt{pdflatex} three times, or
with \texttt{latexmk -pdf}. The text and mathematical fonts are the
same \texttt{newpxtext}/\texttt{newpxmath} family as in the supplied
synthesis, and the five named colors and page geometry are retained.
No font files are distributed. The build checks concern rendering
and cross-references, not the logical correctness of the mathematics.

\clearpage
\phantomsection
\addcontentsline{toc}{section}{References}
\begin{thebibliography}{9}

\bibitem{archive}
\emph{Large cardinals at the inconsistency frontier: a synthesis},
second edition, prepared for Vladimir Reshetnikov, 18 September 2026.
User-supplied source in \nolinkurl{Cardinals3.zip},
\nolinkurl{docs/Large_Cardinals_Synthesis.tex}.
See especially the completeness barrier, ``Two strongly compacts,''
and ``Collapse or destroy.'' Together with
\nolinkurl{docs/Large_Cardinals_Unified_Report.tex} and the
\nolinkurl{Cardinals/} Lean reference development.

\bibitem{cover}
Gabriel Goldberg, reporting joint work with Douglas Blue.
\emph{Consistency beyond the Kunen inconsistency}.
Lecture notes, 1 July 2026.
\url{https://math.berkeley.edu/~goldberg/Slides/CoverExact.pdf}.
Definitions~3.1--3.2, Lemma~3.3, Proposition~3.4, and Problem~5.2.

\bibitem{unique}
Gabriel Goldberg.
\emph{The uniqueness of elementary embeddings}.
The Journal of Symbolic Logic \textbf{89}(4) (2024), 1430--1454.
\doi{10.1017/jsl.2024.60}.
Version checked: \arxiv{2103.13961v2}, revised 2 November 2021;
Proposition~4.2 and Theorems~4.3, 4.9, 4.13.
The corresponding ground theorem in the journal numbering is
Theorem~4.10.

\bibitem{dichotomy}
Gabriel Goldberg.
\emph{Strongly compact cardinals and ordinal definability}.
Journal of Mathematical Logic \textbf{24}(1) (2024), 2250010.
\arxiv{2107.00513}.
Author version checked, 10 August 2021:
\url{https://math.berkeley.edu/~goldberg/Papers/StronglyCompactCardinalsAndOrdinalDefinability.pdf}.
See Lemma~2.1, Proposition~2.3, Lemma~2.7, and Corollary~2.9.

\bibitem{note}
Gabriel Goldberg.
\emph{A note on Woodin's HOD dichotomy}.
\arxiv{2102.05463v4}, 2021.
Proposition~2.2 gives a closely related stationary-seed construction.

\bibitem{abl}
Juan P. Aguilera, Joan Bagaria, and Philipp L\"ucke.
\emph{Large cardinals, structural reflection, and the HOD Conjecture}.
\arxiv{2411.11568v4}, 2025.
Theorem~2.10 gives the basic inner-model regularity phenomenon at
exacting cardinals; the cover-relative version used here is
reconstructed in Section~\ref{sec:barrier}.

\bibitem{kunen}
Kenneth Kunen.
\emph{Elementary embeddings and infinitary combinatorics}.
The Journal of Symbolic Logic \textbf{36}(3) (1971), 407--413.
\doi{10.2307/2269948}.
The local rank-segment form is also stated as Theorem~1.1 in
\cite{cover}.

\end{thebibliography}
\end{document}
